% ============================================================
% Chapter 7: Results
%   Section: Pilot on X-Fact
%     Subsection: Beside the X-Fact Baselines
% Included from chapters/07-results/01-xfact-pilot/xfact-pilot.tex
% ============================================================

\subsection{Beside the X-Fact Baselines}
\label{subsec:pilot-comparison}

Table~\ref{tab:xfact-paper} reproduces the baselines of the X-Fact
paper \cite{gupta2021xfact} from its Tables~3 and~4, read from the
published version. Every figure is a macro-averaged $F_1$ in percent,
the mean of four fine-tuning runs (standard deviation in
parentheses), over X-Fact's seven labels. The paper reports neither
accuracy nor any per-language figure. Its models are multilingual
BERT fine-tuned on 19,079 training claims; ``Meta'' adds the
language, site, claimant and both dates as text, and the evidence
models (Attn-EA) read the snippets of the top five Google results for
each claim, collected when the dataset was built. There is no uniform
random baseline in the paper: ``Attn-EA (Random)'' is the evidence
model given random snippets, and the majority baseline always answers
\emph{false}.

\begin{table}[ht]
\centering
\small
\begin{tabular}{@{}lccc@{}}
\toprule
\textbf{Model (X-Fact paper)} & $\boldsymbol{\alpha_1}$ \textbf{in-domain} & $\boldsymbol{\alpha_2}$ \textbf{out-of-domain} & $\boldsymbol{\alpha_3}$ \textbf{zero-shot} \\
\midrule
Majority & 6.9 & 10.6 & 7.6 \\
Claim-Only & 38.2 (0.9) & 16.2 (0.9) & 14.7 (0.6) \\
Claim-Only + Meta & 39.4 (0.9) & 15.4 (0.8) & 16.7 (1.1) \\
Attn-EA (Random) & 37.5 (0.8) & 16.3 (0.5) & 14.9 (1.2) \\
Attn-EA & 38.9 (0.2) & 15.7 (0.1) & 16.5 (0.7) \\
Attn-EA + Meta & \textbf{41.9} (1.2) & 15.4 (1.5) & 16.0 (0.3) \\
\midrule
\multicolumn{4}{@{}l}{\emph{Training augmented with 12,311 English PolitiFact claims (Table~4)}} \\
Claim-Only + Meta & 37.1 (2.7) & 14.5 (0.5) & 14.4 (0.3) \\
Attn-EA + Meta & 38.0 (4.5) & 14.7 (2.6) & 14.3 (1.9) \\
\bottomrule
\end{tabular}
\caption{The X-Fact baselines as published (Gupta and Srikumar,
Tables~3 and~4): macro-$F_1$ in percent over seven labels, mean
(standard deviation) of four runs. $\alpha_1$: 3,826 claims in the
training languages and sources, Spanish (Chequeado) among them;
$\alpha_2$: 2,368 claims from unseen sources of four training
languages; $\alpha_3$: 3,381 claims in twelve unseen languages.}
\label{tab:xfact-paper}
\end{table}

Table~\ref{tab:xfact-beside} sets the system's macro-$F_1$ beside the
paper's closest rows. The Spanish test split is the only row of ours
drawn the way $\alpha_1$ is, and it covers one of $\alpha_1$'s twelve
languages; the pilot's English half has no counterpart in the paper at
all, since English is not among X-Fact's evaluation languages.

\begin{table}[ht]
\centering
\small
\begin{tabularx}{\linewidth}{@{}Xrr@{}}
\toprule
\textbf{System and set} & \textbf{Labels} & \textbf{Macro-$F_1$ ($\times$100)} \\
\midrule
\multicolumn{3}{@{}l}{\emph{X-Fact paper, supervised mBERT, mean of 4 runs}} \\
Majority, $\alpha_1$ & 7 & 6.9 \\
Claim-Only, $\alpha_1$ / $\alpha_3$ & 7 & 38.2 / 14.7 \\
Attn-EA (claim + snippets), $\alpha_1$ / $\alpha_3$ & 7 & 38.9 / 16.5 \\
Attn-EA + Meta (best), $\alpha_1$ / $\alpha_3$ & 7 & 41.9 / 16.0 \\
\midrule
\multicolumn{3}{@{}l}{\emph{This work, zero-shot \texttt{llama3.2:3b}, one run, 95\% interval}} \\
Spanish test split ($\alpha_1$, Spanish only) & 5 & 11.1 [7.9, 14.2] \\
Pilot, Spanish & 5 & 13.0 [6.1, 19.6] \\
Pilot, English (not in X-Fact's evaluation) & 3 & 35.8 [15.7, 55.4] \\
Closed-book (no evidence), Spanish test split & 5 & 8.4 [4.9, 12.7] \\
Majority, Spanish test split & 5 & 9.3 \\
\bottomrule
\end{tabularx}
\caption{The system beside the X-Fact baselines. The two halves differ
in label set, training, evidence, sample and year of search; the list
that follows says how, and none of the differences is small.}
\label{tab:xfact-beside}
\end{table}

On the one row drawn as $\alpha_1$ is, the system's macro-$F_1$ of 11.1
is 27 points below the paper's claim-only model and 31 below its best,
above its majority baseline, and below what a uniform random guess is
expected to score on the same claims (about 19). Read through the
differences listed below, this says that the zero-shot pipeline does
not approach a supervised model on X-Fact's own distribution; it does
not say by how much the two would differ on equal terms.

The two halves of Table~\ref{tab:xfact-beside} are not directly
comparable, for reasons that each move the figure by more than its
interval:

\begin{itemize}
    \item \textbf{Labels.} X-Fact scores seven labels; this work scores
    five, with \emph{mostly false} merged into \texttt{FALSE},
    \emph{mostly true} mapped to \texttt{PARTIALLY\_TRUE} and
    \emph{other} dropped. A merged class is easier to hit, and
    macro-$F_1$ averages over fewer classes.
    \item \textbf{English.} English rows exist only in X-Fact's
    training data, keep PolitiFact's own six-point scale (no
    \emph{partly true/misleading}, no \emph{complicated}) and map
    PolitiFact's \emph{half true} to \texttt{PARTIALLY\_TRUE}, whereas
    X-Fact's master mapping sends \emph{half true} to \emph{partly
    true/misleading}, which this work maps to \texttt{MISLEADING}.
    \item \textbf{Sample.} The pilot is stratified, eight claims per
    language and label, so its labels are balanced by construction;
    X-Fact's test sets are not. Only the Spanish test split is drawn
    as $\alpha_1$ is.
    \item \textbf{Size.} The paper's figures are means over four runs
    on thousands of claims, and their standard deviations (0.1 to 4.5
    points) measure seeds, not sampling; this work's are single runs
    over 64 or 165 claims, with intervals 13 to 40 points wide.
    \item \textbf{Supervision.} mBERT was fine-tuned on 19,079 X-Fact
    claims and learned their label distribution and the claim-only
    cues the paper itself documents (its claim-only model beats the
    majority by 31 points on $\alpha_1$). The model here never saw an
    X-Fact label, which makes $\alpha_3$, where the supervised models
    also face unseen data, the fairer reference despite not containing
    Spanish.
    \item \textbf{Evidence and its date.} The paper's evidence is the
    snippets of the top five Google results, collected when the dataset
    was built; this work's is the open web of 13 October 2026, through
    the engines that answered SearXNG that day (Google and Brave did
    not), fetched in full and filtered by the pertinence gate.
    Section~\ref{subsec:pilot-time} is about what that changes.
    \item \textbf{Metadata and context.} The best paper model reads the
    claimant, the site and both dates; the pipeline reads the claim
    text alone, and Chequeado's claims are quotes stripped of speaker,
    place and moment.
\end{itemize}
